% GENERATED FILE. Edit canonical lessons, then run npm run generate:books.
% canonical-source-hash: fnv1a64:9372026223ee689e
% canonical-lessons: GU-C223-tabiyat, GU-C223-sharir, GU-C223-dukhvu, GU-C223-shardi, GU-C223-chikvu

\chapter{How are you feeling}
\label{ch:gu-how-are-you-feeling}
\hlchaptermodality{223}

\begin{chapteropening}
\textbf{By the end of this chapter:} \emph{I can ask about someone's health and say where it hurts.}

Complete the last lesson of chapter 223: I can ask about someone's health and say where it hurts.
\end{chapteropening}

\section[\texorpdfstring{\gu{તબિયત} (tabiyat)}{તબિયત (tabiyat)}]{\gu{તબિયત} (tabiyat) — health}

\label{lesson:GU-C223-tabiyat}



\begin{quote}
Before the new one: say the Gujarati for high tide, then the Gujarati for low tide.
\end{quote}

\subsection*{You'll want to know: \gu{તબિયત}}
\textbf{\gu{તબિયત}} — \emph{tabiyat} — \textquotedblleft{}health\textquotedblright{}.

\begin{quote}
\textbf{\gu{શું} \gu{તમારી} \gu{તબિયત} \gu{સારી} \gu{છે}?} — \emph{śũ tamārī tabiyat sārī che?} — \textquotedblleft{}Are you keeping well?\textquotedblright{}
\end{quote}

Word for word, is your health good.

\begin{cousinweb}
\textbf{\gu{તબિયત}} (\emph{tabiyat}) is a loan from Arabic \emph{ṭabīʿa}, \textquotedblleft{}nature, temperament\textquotedblright{}, which reached Gujarati through Persian (EtymDB lists the Persian word). Its Arabic root is \emph{ṭ-b-ʿ} (Lane).
\end{cousinweb}

\begin{grammarlens}[title={What you've built}]
One more word for everyday life at A2.
\end{grammarlens}

\subsection*{Your turn}
\begin{itemize}
\raggedright
  \item \emph{Say it:} \emph{tabiyat}
  \item \emph{Say it:} \emph{tabiyat}, once more
  \item \emph{Say it:} \emph{śũ tamārī tabiyat sārī che?} — Are you keeping well?
  \item \emph{Recall:} say the Gujarati for high tide, then the Gujarati for low tide, then say \emph{tabiyat} again
\end{itemize}

\begin{tcolorbox}[breakable,colback=teal!4,colframe=teal!35!black,title={Before you move on}]
What does \gu{તબિયત} mean? (\textquotedblleft{}Health\textquotedblright{}.) Say it once more.
\end{tcolorbox}
\section[\texorpdfstring{\gu{શરીર} (śarīr)}{શરીર (śarīr)}]{\gu{શરીર} (śarīr) — body}

\label{lesson:GU-C223-sharir}



\begin{quote}
Before the new one: say the Gujarati for low tide, then the Gujarati for health.
\end{quote}

\subsection*{You'll want to know: \gu{શરીર}}
\textbf{\gu{શરીર}} — \emph{śarīr} — \textquotedblleft{}body\textquotedblright{}.

\begin{quote}
\textbf{\gu{મારું} \gu{શરીર} \gu{નબળું} \gu{નથી}.} — \emph{mārũ śarīr nabaḷũ nathī} — \textquotedblleft{}My body is not weak.\textquotedblright{}
\end{quote}

\begin{cousinweb}
\textbf{\gu{શરીર}} (\emph{śarīr}) is Sanskrit \emph{śarīra}, \textquotedblleft{}body\textquotedblright{}, taken over whole as a learned loan, a tatsama. CDIAL lists the inherited Old Gujarati \emph{saïra}, \textquotedblleft{}body\textquotedblright{}, beside it (CDIAL 12335).
\end{cousinweb}

\begin{grammarlens}[title={What you've built}]
One more word for everyday life at A2.
\end{grammarlens}

\subsection*{Your turn}
\begin{itemize}
\raggedright
  \item \emph{Say it:} \emph{śarīr}
  \item \emph{Say it:} \emph{śarīr}, once more
  \item \emph{Say it:} \emph{mārũ śarīr nabaḷũ nathī} — My body is not weak.
  \item \emph{Recall:} say the Gujarati for low tide, then the Gujarati for health, then say \emph{śarīr} again
\end{itemize}

\begin{tcolorbox}[breakable,colback=teal!4,colframe=teal!35!black,title={Before you move on}]
What does \gu{શરીર} mean? (\textquotedblleft{}Body\textquotedblright{}.) Say it once more.
\end{tcolorbox}
\section[\texorpdfstring{\gu{દુખવું} (dukhvũ)}{દુખવું (dukhvũ)}]{\gu{દુખવું} (dukhvũ) — to ache, to hurt}

\label{lesson:GU-C223-dukhvu}



\begin{quote}
Before the new one: say the Gujarati for health, then the Gujarati for body.
\end{quote}

\subsection*{You'll want to know: \gu{દુખવું}}
\textbf{\gu{દુખવું}} — \emph{dukhvũ} — \textquotedblleft{}to ache, to hurt\textquotedblright{}.

\begin{quote}
\textbf{\gu{શું} \gu{તમારું} \gu{માથું} \gu{દુખે} \gu{છે}?} — \emph{śũ tamārũ māthũ dukhe che?} — \textquotedblleft{}Does your head ache?\textquotedblright{}
\end{quote}

The noun for the pain is \textbf{\gu{દુખાવો}}.

\begin{cousinweb}
\textbf{\gu{દુખવું}} (\emph{dukhvũ}) is a tadbhava, a word inherited from Sanskrit and reshaped by sound change. It comes from \emph{duḥkhati}, \textquotedblleft{}pains\textquotedblright{}, built on \emph{duḥkha}, \textquotedblleft{}pain\textquotedblright{} (CDIAL 6376).
\end{cousinweb}

\begin{grammarlens}[title={What you've built}]
One more word for everyday life at A2.
\end{grammarlens}

\subsection*{Your turn}
\begin{itemize}
\raggedright
  \item \emph{Say it:} \emph{dukhvũ}
  \item \emph{Say it:} \emph{dukhvũ}, once more
  \item \emph{Say it:} \emph{śũ tamārũ māthũ dukhe che?} — Does your head ache?
  \item \emph{Recall:} say the Gujarati for health, then the Gujarati for body, then say \emph{dukhvũ} again
\end{itemize}

\begin{tcolorbox}[breakable,colback=teal!4,colframe=teal!35!black,title={Before you move on}]
What does \gu{દુખવું} mean? (\textquotedblleft{}To ache, to hurt\textquotedblright{}.) Say it once more.
\end{tcolorbox}
\section[\texorpdfstring{\gu{શરદી} (śardī)}{શરદી (śardī)}]{\gu{શરદી} (śardī) — common cold, a chill}

\label{lesson:GU-C223-shardi}



\begin{quote}
Before the new one: say the Gujarati for body, then the Gujarati for to ache.
\end{quote}

\subsection*{You'll want to know: \gu{શરદી}}
\textbf{\gu{શરદી}} — \emph{śardī} — \textquotedblleft{}common cold, a chill\textquotedblright{}.

\begin{quote}
\textbf{\gu{મને} \gu{શરદી} \gu{નથી}.} — \emph{mane śardī nathī} — \textquotedblleft{}I don't have a cold.\textquotedblright{}
\end{quote}

A feminine noun.

\begin{cousinweb}
\textbf{\gu{શરદી}} (\emph{śardī}) is a loan, usually traced to Persian \emph{sardī}, \textquotedblleft{}coldness\textquotedblright{}, from \emph{sard}, \textquotedblleft{}cold\textquotedblright{}; EtymDB lists the Persian word. CDIAL has no entry for it.
\end{cousinweb}

\begin{grammarlens}[title={What you've built}]
One more word for everyday life at A2.
\end{grammarlens}

\subsection*{Your turn}
\begin{itemize}
\raggedright
  \item \emph{Say it:} \emph{śardī}
  \item \emph{Say it:} \emph{śardī}, once more
  \item \emph{Say it:} \emph{mane śardī nathī} — I don't have a cold.
  \item \emph{Recall:} say the Gujarati for body, then the Gujarati for to ache, then say \emph{śardī} again
\end{itemize}

\begin{tcolorbox}[breakable,colback=teal!4,colframe=teal!35!black,title={Before you move on}]
What does \gu{શરદી} mean? (\textquotedblleft{}Common cold, a chill\textquotedblright{}.) Say it once more.
\end{tcolorbox}
\section[\texorpdfstring{\gu{છીંકવું} (chī̃kvũ)}{છીંકવું (chī̃kvũ)}]{\gu{છીંકવું} (chī̃kvũ) — to sneeze}

\label{lesson:GU-C223-chikvu}



\begin{quote}
Before the new one: say the Gujarati for to ache, then the Gujarati for common cold.
\end{quote}

\subsection*{You'll want to know: \gu{છીંકવું}}
\textbf{\gu{છીંકવું}} — \emph{chī̃kvũ} — \textquotedblleft{}to sneeze\textquotedblright{}.

\begin{quote}
\textbf{\gu{હું} \gu{આજે} \gu{નથી} \gu{છીંકતો}.} — \emph{hũ āje nathī chī̃kto} — \textquotedblleft{}I am not sneezing today.\textquotedblright{}
\end{quote}

The sneeze is \textbf{\gu{છીંક}}, \emph{chī̃k}.

\begin{cousinweb}
\textbf{\gu{છીંકવું}} (\emph{chī̃kvũ}) is inherited, built on an imitative word for the sound. CDIAL lists Gujarati \emph{chī̃k} and \emph{chī̃kvũ} under Sanskrit \emph{chikkā}, \textquotedblleft{}sneeze\textquotedblright{} (CDIAL 5032).
\end{cousinweb}

\begin{grammarlens}[title={What you've built}]
One more word for everyday life at A2.
\end{grammarlens}

\subsection*{Your turn}
\begin{itemize}
\raggedright
  \item \emph{Say it:} \emph{chī̃kvũ}
  \item \emph{Say it:} \emph{chī̃kvũ}, once more
  \item \emph{Say it:} \emph{hũ āje nathī chī̃kto} — I am not sneezing today.
  \item \emph{Recall:} say the Gujarati for to ache, then the Gujarati for common cold, then say \emph{chī̃kvũ} again
\end{itemize}

\begin{tcolorbox}[breakable,colback=teal!4,colframe=teal!35!black,title={Before you move on}]
What does \gu{છીંકવું} mean? (\textquotedblleft{}To sneeze\textquotedblright{}.) Say it once more.
\end{tcolorbox}
